\chapter{Materiais e métodos}\label{cap:metodologia}

\section{Natureza e abordagem da pesquisa}

Este trabalho caracteriza-se como pesquisa aplicada, de natureza tecnológica,
orientada à análise de um artefato de \emph{software} de código aberto e à sua
avaliação em contexto de uso real. A abordagem combina elementos
quantitativos (contagens, taxas, resultados de verificação) e qualitativos
(análise de decisões de projeto, limitações e implicações). A avaliação adota
o formato de \textbf{estudo de caso instrumental}: um projeto de placa é
utilizado como instrumento para examinar as capacidades e os limites do
servidor, e não como amostra representativa de todos os projetos de PCB. O
documento segue as normas ABNT NBR 14724~\cite{abnt14724},
NBR 6023~\cite{abnt6023} e NBR 10520~\cite{abnt10520} para estrutura,
referências e citações.

\section{Percurso metodológico}

O percurso foi organizado em cinco etapas:

\begin{enumerate}[leftmargin=1.4em,itemsep=0.25em]
  \item \textbf{E1 --- Auditoria do domínio e do código}: leitura da
        documentação do servidor, do inventário de ferramentas, da arquitetura
        e da suíte de testes; identificação dos pontos de integração com o
        KiCad e com APIs externas.
  \item \textbf{E2 --- Definição de requisitos}: consolidação dos requisitos
        R1 a R6 (Seção~\ref{sec:req-monografia}) a partir das lacunas
        identificadas na fundamentação.
  \item \textbf{E3 --- Análise da implementação}: leitura do código existente
        no \emph{workspace}, incluindo a arquitetura em duas camadas, o
        protocolo interno e os mecanismos de robustez documentados no código
        e no histórico do projeto.
  \item \textbf{E4 --- Verificação automatizada}: execução das suítes
        TypeScript e Python com o ambiente completo do projeto, em
        10/09/2026, com registro de resultados em arquivo de evidência.
  \item \textbf{E5 --- Estudo de caso}: caracterização da placa seguidora de
        linhas, verificação de regras de projeto com \texttt{kicad-cli},
        geração de renderizações 3D e inspeção dos artefatos de fabricação.
\end{enumerate}

\section{Requisitos do sistema}\label{sec:req-monografia}

O Quadro~\ref{qua:requisitos} apresenta os requisitos que orientaram a
avaliação, com a forma de atendimento observada e as evidências
correspondentes.

\begin{quadro}[htbp]
\centering
\small
\caption{Requisitos, atendimento e evidências.}
\label{qua:requisitos}
\begin{tabular}{@{}p{0.7cm}p{3.4cm}p{5.2cm}p{4.6cm}@{}}
\toprule
\textbf{ID} & \textbf{Requisito} & \textbf{Atendimento} & \textbf{Evidência} \\
\midrule
R1 & Cobertura do fluxo de PCB & 233 ferramentas em 24 módulos funcionais (projeto a fabricação) & Inventário do servidor (06/09/2026) \\
R2 & Descoberta de operações & Registro de 16 categorias, lista de essenciais e busca por palavra-chave & 173 ferramentas indexadas; Capítulo~\ref{cap:catalogo} \\
R3 & Robustez sob falhas & Correlação de requisições, descarte de respostas tardias, tempos-limite por comando & Testes TypeScript de correlação e de inicialização \\
R4 & Integridade dos arquivos & Salvamento explícito, sessão fixada por \emph{backend}, validação estrutural de esquemáticos & Suíte Python; seção de decisões de engenharia \\
R5 & Portabilidade & Node.js e Python portáveis; configuração por variáveis de ambiente & Matriz de integração contínua (4 sistemas para TypeScript) \\
R6 & Verificabilidade & Testes em duas linguagens e inventário gerado como portão de qualidade & 84 casos TypeScript; 2.367 casos Python; \texttt{docs:tools:check} \\
\bottomrule
\end{tabular}
\end{quadro}

\section{Métricas e instrumentos}

As métricas adotadas e seus instrumentos são descritos no
Quadro~\ref{qua:metricas}. Todas são reprodutíveis a partir do repositório.

\begin{quadro}[htbp]
\centering
\small
\caption{Métricas de avaliação e instrumentos de coleta.}
\label{qua:metricas}
\begin{tabular}{@{}p{3.8cm}p{6.2cm}p{3.9cm}@{}}
\toprule
\textbf{Métrica} & \textbf{Definição} & \textbf{Instrumento / fonte} \\
\midrule
Cobertura de descoberta & razão entre ferramentas indexadas e registradas & inventário gerado por \texttt{scripts/generate-tool-inventory.mjs} \\
Tamanho de implementação & linhas de código por camada e por módulo & contagem direta em \texttt{src/} e \texttt{python/} \\
Casos de teste & casos executados, aprovados, com falha e ignorados & Vitest e pytest (logs em \texttt{evidencias/}) \\
Conformidade de projeto & violações de DRC por tipo e severidade & \texttt{kicad-cli pcb drc --format json --severity-all} \\
Conectividade & itens não conectados e paridade esquemático--placa & relatório de DRC \\
Artefatos de fabricação & presença e tipo de arquivos de produção & inspeção do diretório \texttt{production/} \\
\bottomrule
\end{tabular}
\end{quadro}

\section{Ambiente de validação}

A validação foi executada em estação de trabalho com Ubuntu 24.04.4 LTS
(\emph{kernel} 7.0.0-31), KiCad 9.0.7 instalado por pacote \emph{snap}
(revisão 22), Node.js~\cite{nodejs} v22.23.2, npm 12.0.2 e
Python~\cite{python} 3.12.3 em ambiente virtual
dedicado. As dependências declaradas pelo projeto (produção e
desenvolvimento) foram instaladas previamente à execução, incluindo as
bibliotecas \texttt{kicad-python}, \texttt{kicad-skip}, \texttt{sexpdata},
\texttt{pymupdf} e \texttt{cairosvg}. O \emph{backend} de execução foi
selecionado automaticamente pela fábrica de \emph{backends}, recaindo sobre a
API \texttt{pcbnew} (SWIG), uma vez que o servidor IPC do KiCad não estava
habilitado no sistema utilizado (detalhes na Seção~\ref{sec:limitacoes}).

Os registros de execução das suítes de teste e o relatório de DRC estão
preservados em \texttt{evidencias/}, com os nomes
\texttt{testes-pytest-2026-09-10.log}, \texttt{testes-vitest-2026-09-10.log}
e \texttt{drc-seguidor-de-linhas-2026-09-10.json}.

\section{Estudo de caso e procedimentos}

A placa utilizada como estudo de caso é o projeto
\texttt{seguidor\_de\_linhas\_4camada2-2026-1}, projetado pelo autor com o
auxílio do servidor no mesmo \emph{workspace}. O projeto destina-se a um robô seguidor de
linhas, com os seguintes blocos funcionais identificados no esquemático e na
lista de materiais: sensoriamento óptico reflexivo; conversão de energia com
rastreamento de ponto de máxima potência (MPPT); acionamento de motores por
pontes H; alimentação com reguladores lineares (LDO); interface USB-C;
medição e sinalização de tensão da bateria; e alerta sonoro.

Os procedimentos executados foram:

\begin{enumerate}[leftmargin=1.4em,itemsep=0.25em]
  \item \textbf{P1 --- Caracterização do projeto}: contagem de
        \emph{footprints} e redes no arquivo \texttt{.kicad\_pcb}; verificação
        da espessura e do número de camadas; contagem de folhas do
        esquemático; inspeção da lista de materiais.
  \item \textbf{P2 --- Verificação de regras}: execução de
        \texttt{kicad-cli pcb drc} com formato JSON e severidade completa,
        exportando o relatório para arquivo de evidência; sumarização por
        tipo e severidade.
  \item \textbf{P3 --- Geração de figuras}: renderizações 3D (faces superior
        e inferior) com \texttt{kicad-cli pcb render} e exportação das folhas
        do esquemático com \texttt{kicad-cli sch export pdf}.
  \item \textbf{P4 --- Execução das suítes}: \texttt{npm run test:ts}
        (Vitest~\cite{vitest}) e \texttt{python -m pytest tests/}
        (pytest~\cite{pytest}), com o ambiente completo descrito, registrando
        os resultados em log.
\end{enumerate}

Os comandos exatos estão reproduzidos no Apêndice~\ref{ap:repro}.

\section{Matriz de rastreabilidade}

O Quadro~\ref{qua:rastreabilidade} relaciona objetivos específicos, métodos,
evidências e seções de resultado, conforme exigência de rastreabilidade
metodológica.

\begin{quadro}[htbp]
\centering
\small
\caption{Matriz de rastreabilidade objetivo--método--evidência.}
\label{qua:rastreabilidade}
\begin{tabular}{@{}p{4.6cm}p{3.4cm}p{3.8cm}p{2.6cm}@{}}
\toprule
\textbf{Objetivo específico} & \textbf{Método} & \textbf{Evidência} & \textbf{Seção} \\
\midrule
Arquitetura analisada & E1/E3 & Capítulo~\ref{cap:arquitetura} & 5 \\
Catálogo caracterizado & E1/E3 & Inventário; Capítulo~\ref{cap:catalogo} & 6 \\
Descoberta e recursos & E4 & Métricas do Capítulo~\ref{cap:resultados} & 8 \\
Testes e integração contínua & E4 & Logs; matriz de CI & 8 \\
Estudo de caso & E5 & DRC; figuras; artefatos & 7 \\
Limitações e continuidade & E5 & Capítulo~\ref{cap:discussao} & 9/10 \\
\bottomrule
\end{tabular}
\end{quadro}

\section{Síntese do capítulo}

O capítulo fixou o desenho metodológico, os requisitos, as métricas e o
ambiente. Os capítulos seguintes apresentam o artefato (arquitetura e
catálogo) e, em seguida, os resultados obtidos nos procedimentos P1 a P4.
